% titel: Blatt 3, A6: Verdünnen (Massen- und Stoffmengenkonzentration)
% skript: Mortimer Kap. 4.6 (Konzentrationen, Verdünnen)

% Teil I: Werkzeuge, kein Lösungsweg.

Eine Stammlösung wird mit Wasser auf ein Zielvolumen aufgefüllt. Dabei bleibt
die Stoffmenge des gelösten Stoffes erhalten; es ändert sich nur die
Konzentration.

\begin{werkzeugkasten}[Konzentrationen und Verdünnen]
Stoffmengenkonzentration und Massenkonzentration:
\begin{equation}
  c = \frac{n}{V}, \qquad
  b = \frac{m_{\text{Stoff}}}{V}
  \qquad (b \text{ heisst hier Massenkonzentration}).
\end{equation}
Übergang zwischen beiden über die molare Masse:
$c = b / M$ (mit $m = n\cdot M$).
Beim Verdünnen ist $n$ konstant: aus einer Stammlösung der
Konzentration $c_0$ entnimmt man das Volumen $V_0$ und füllt auf das Zielvolumen
$V$ mit Zielkonzentration $c$ auf:
\begin{equation}
  n = c_0 V_0 = c\,V \;\Longrightarrow\;
  V_0 = \frac{c\,V}{c_0}.
\end{equation}
\textit{Skript:} Mortimer Kap. 4 (Konzentrationen, Verdünnen), Übung 3/6.
\end{werkzeugkasten}

\begin{achtung}[Darauf achten]
\begin{itemize}
  \item Beide Konzentrationen in \emph{dieselbe} Einheit bringen, bevor man
        sie vergleicht oder $c = b/M$ rechnet (hier: Massenkonzentration
        $750\ \mathrm{mg/L}$ in mol/L umrechnen).
  \item Vorspräfixe: mmol/L $= 10^{-3}\ \mathrm{mol/L}$, mg/L
        $= 10^{-3}\ \mathrm{g/L}$; Molmasse in g/mol, Masse in g.
  \item Die Stoffmenge $n = c\cdot V$ wandert beim Verdünnen unverändert von
        der entnommenen Stammlösung in das Zielgefäss.
  \item Ergebnis plausibel: entnommenes Volumen $V_0 < V$, sonst wäre es kein
        Verdünnen.
\end{itemize}
\end{achtung}
